% ============================================================
% PHẦN 3.5 — ĐỊNH VỊ VÀ THÔNG ĐIỆP TRUYỀN THÔNG
% ------------------------------------------------------------
% Thuộc: PA3-02
% Người phụ trách chính: Nguyễn Anh Thư – CEO/Project Lead
% Phối hợp: Trần Minh Quang
%
% CÂU HỎI CẦN TRẢ LỜI:
%   - 1Stream đang giải quyết vấn đề cụ thể nào cho trung tâm ngoại ngữ?
%   - Giá trị chính là: Giảm phụ thuộc vào host? Mở rộng thời lượng phát?
%     Chuẩn hóa thông tin tuyển sinh? Thu lead? Hay giảm khối lượng công việc lặp lại?
%   - Khách hàng nên hiểu 1Stream khác gì so với:
%       Agency livestream / Nhân viên nội bộ / Công cụ tạo video AI /
%       Công cụ phát livestream / Chatbot?
%   - Có bằng chứng nào hỗ trợ thông điệp?
%   - Những lợi ích nào hiện mới chỉ là giả thuyết?
%
% OUTPUT BẮT BUỘC:
%   1. Tuyên bố định vị (Positioning Statement):
%      "Dành cho [khách hàng mục tiêu], đang gặp [vấn đề],
%       1Stream là [loại giải pháp], giúp [kết quả],
%       bằng cách [cơ chế khác biệt]."
%
%   2. Thông điệp truyền thông CHÍNH (1 câu)
%
%   3. Ba thông điệp HỖ TRỢ, ví dụ:
%      - Giảm công việc lặp lại
%      - Mở rộng thời lượng livestream
%      - Chuẩn hóa tư vấn và thu lead
%
%   4. Bảng thông điệp theo từng người trong đơn vị mua:
%      - Chủ trung tâm
%      - Trưởng marketing
%      - Nhân viên vận hành
%      - Tư vấn tuyển sinh
%
% TIÊU CHÍ HOÀN THÀNH:
%   - Không dùng thông điệp "nhiều tính năng hơn đối thủ"
%   - Không khẳng định hiệu quả doanh thu khi chưa có pilot
%   - Phân biệt rõ: Lời hứa / Cơ chế tạo giá trị / Bằng chứng đã có / Giả định cần kiểm chứng
% ============================================================

\subsection{Định vị và thông điệp truyền thông}

% TODO (Nguyễn Anh Thư): Viết nội dung mục 3.5
